% -*- coding: utf-8 -*-
% 线代基础知识.tex
% 线代基础知识
\documentclass[UTF8]{ctexart}
\usepackage{amsmath}
\usepackage{xcolor}
\usepackage{geometry}


\geometry{a4paper,scale=0.8}
\title{线性代数基础知识}
\author{埃及猪肉}
\date{\today}

\begin{document}
\maketitle
\tableofcontents

\section{矩阵}

\subsection{矩阵的运算}

\paragraph{加法和数乘}
这两种运算和代数运算一样
\begin{itemize}
    \item 交换律：
        \[A + B = B + A\]
    \item 结合律：
        \[(A + B) + C = A + (B + C)\]
    \item 分配律：
        \begin{align*}
        k(A + B) &= kA + kB \\
        (k + l)A &= kA + lA
        \end{align*}
\end{itemize}

\paragraph{矩阵乘法}
\begin{itemize}
    \item 结合律：
        \[
        (A_{m \times s} B_{s \times r}) C_{r \times n} = A_{m \times s} (B_{s \times r} C_{r \times n})
        \]
    \item 分配律：
        \[
        A_{m \times s} (B_{s \times n} + C_{s \times n}) = A_{m \times s} B_{s \times n} + A_{m \times s} C_{s \times n}
        \]
        \[
        (A_{m \times s} + B_{m \times s}) C_{s \times n} = A_{m \times s} C_{s \times n} + B_{m \times s} C_{s \times n}
        \]
\end{itemize}

{\color{red}
\paragraph{关键：}矩阵乘法没有交换律，其他都一样。
\[
AB \neq BA.
\]
}

（但$(kA_{m \times s}) B_{s \times n} = A_{m \times s} (kB_{s \times n}) = k (A_{m \times s} B_{s \times n})$常数仍能使用交换律，但不能改变矩阵顺序。）




\end{document}
